\cotefile{97}
\scanpage{5}

Classes de Chern de la représentation régulière et ses multiples.

$G$ groupe cyclique d'ordre $n$,
\[
H^*(G,\mathbb Z)^+ \simeq
\underline{\mathrm{Sym}}_{\mathbb Z_{\! /n\mathbb Z}}
  (\overset{\text{\scriptsize\inserted{dual sur $\mathbb Z_{\! /n\mathbb Z}$}}}{\check G})^+
\simeq \mathbb Z_{\! /n\mathbb Z}[\xi]^+
\]
où $\xi$ est un générateur de
$H^2(G,\mathbb Z)\simeq
\underset{\substack{\rotatebox{90}{$\simeq$}\\ \check G}}{
\mathrm H^{\deleted{\unsure{2|*}}}\mathrm{om}(G,\mathbb Q/\mathbb Z)}$

La représentation régulière \inserted{(sur $\mathbb C$)} \deleted{est} correspond à la somme des représentations de degré $1$ \unsure{associés|associées} aux \unsure{différents} $\xi\in\check G$.

On a donc
\[
c(r_G)=\prod_{\xi\in\check G}(1+\xi)
=\prod_{\substack{\deleted{{\unsure{k\in\mathbb Z|k=0}}^{\unsure{n-1|n}}}\\ \lambda\in\mathbb Z_{\! /n\mathbb Z}}}(1+\lambda\xi_0)
\deleted{=\unsure{\sum_{i=0}^{n}c_i\xi_0^i\;(\text{dans }\mathbb Z_{\! /n\mathbb Z}[\xi])}}
\]

Supposons $n=p$ premier. On trouve
\[
\left\{\begin{aligned}
c_p(r_G)&=0\\
c_{p-1}(r_G)&=\prod_{\xi\in\check G-\{0\}}\xi
=\left(\prod_{\lambda\in\mathbb F_p^{\unsure{+|*}}}\lambda\right)\xi_0^{p-1}
=\xi_0^{p-1}\quad \text{(indépendant du choix de $\xi_0\in\check G$ !)}
\end{aligned}\right.
\]

D'où
\[
\left\{\begin{aligned}
c_{\deleted{\unsure{p}}i}(Nr_G)&=0\quad\text{si }i>N(p-1)\\
c_{N(p-1)}(Nr_G)&=\xi^{N(p-1)}\qquad(\xi\in\check G)
\end{aligned}\right.
\]

Soit maintenant $G$ un groupe fini\inserted{ d'ordre $n$}, d'ordre divisible par $p$, \deleted{Alors on \unsure{trouve|tous}} et soit $H$ un sous-groupe d'ordre $p$. Considérons $n'=n/p$, alors $\deleted{(Nr_G)\vert}r_G\vert H=n'r_H$ donc $(Nr_G)\vert H=(Nn')r_H$, d'où $\bigl(c_i(nr_G)\bigr)\vert H=c_i(Nn'r_H)$, et par suite
\[
c_{Nn'(p-1)}(r_G)\vert H\ne0,\quad\text{à fortiori}\quad c_{Nn'(p-1)}(r_G)\ne0
\]
Donc
\[
\boxed{H^{\deleted{\unsure{*|2}}N\unsure{K|k}}(G)(p)\ne0\quad\text{pour tout entier }N,\quad\text{où}\quad k=2n'(p-1)}
\]
% archive sheet number: 4
\scanpage{7}
\mspage{\textcircled{1}}
% Faint mirror-image typescript shows through from the other side.
\diagram{\begin{tikzcd}[ampersand replacement=\&,column sep=large,row sep=small]
\unsure{\widetilde S|{\widetilde{\mathbb S}}}\arrow[r] \& \unsure{S|\mathbb S}\\
\widetilde G=\mathrm{Sl}(2,\mathbb R)\arrow[r]\arrow[d,phantom,"\unsure{\cup|\downarrow|\vert}" description] \&
G=\mathrm{GP}(2,\mathbb R)=\widetilde G/\deleted{(\unsure{\pm1})}(+1,-1)
\arrow[d,phantom,"\cup" description]\\
\widetilde\Gamma\arrow[r,"\sim"] \& \Gamma
\end{tikzcd}}

\diagram{\begin{tikzcd}[ampersand replacement=\&,row sep=small]
\& V\\
D\arrow[d]\&\\
X\&
\end{tikzcd}}

$T_X$, $T_D$ avec op. de $\Gamma\deleted{\unsure{=\widetilde\Gamma}}$, induit par $T_D$ avec op. de $G$ déterminé par op. de $S=G_x$ sur $T_D(x)$

\inserted{$\widetilde\Gamma$ induit par $T_D$ ``avec'' op. de $\widetilde G$, \unsure{déterminé|déduit} par op. $\chi$ de $\widetilde S=\widetilde G_x\longrightarrow S$ sur $T_D(x)$}

$E_0$ défini par repr de $\widetilde\Gamma$ induite par $\widetilde G$ op. sur $V$ déterminé par fibré vectoriel \unsure{$c^t$|$c^{st}$} : $D\times V$ avec op diagonale de $\widetilde\Gamma$ $\gamma.(x,v)=(\gamma x,\gamma v)$, induite par op. analogue de $\widetilde G$ sur $D\times V$, \unsure{définie|déduite} par \unsure{représentation|l'opération|$V$ vue} \unsure{canon|comme} $\chi_0$ de $\widetilde S$ sur $V$.

On a
\[
\boxed{\chi=\chi_0^{(2)}}
\]
\[
\begin{aligned}
\text{Caractère de }\chi_0 &: e_1+e_{-1}\\
\text{Caractère de }\chi &: \deleted{\unsure{e_1+e_{-1}}}\;e_2+e_{-2}\\
\chi&=\chi_0^2-2
\end{aligned}
\]

\deleted{\unsure{fais|Lusin}} \uline{N.B.} Pour tte \deleted{\unsure{const.}} représentation linéaire réelle \deleted{(\unsure{réelle de dim. 2|rep. complexe de})}\inserted{(resp complexe)} de $\widetilde S$, grâce à $\widetilde\Gamma\simeq\Gamma\subset\widetilde G$, correspond un fibré vectoriel\inserted{ réel} \deleted{(\unsure{réel|triv.})}\inserted{(resp complexe -- mais pas holomorphe)} sur $X$. Cette \unsure{correspondance|corresp. est un}
% archive sheet number: 6
\scanpage{9}
% Faint reversed typescript shows through.
respecte les opérations tensorielles, l'opération de \unsure{complexification|conjugaison}.

Or $\chi_0$ correspond à \unsure{une|un} fibré vectoriel complexe de rang $1$, et $\chi$ correspond au \uline{carré} du $\unsure{=|\text{---}}$.

\deleted{Par suite} Si $E\inserted{=T_X}$, $E_0$ sont les fibrés vectoriels \deleted{réels}\inserted{complexes} associés, leurs classes de Chern \deleted{\unsure{sont}} $\widetilde c_1$, $c_1$ sont \unsure{reliées|reliés} par
\[
\widetilde c_1=2c_1
\]
d'autre part ce sont aussi les classes \deleted{de Chern}\inserted{caractéristiques} des fibrés vectoriels associés. Comme \inserted{pour} $E=T_X$, on \unsure{trouve} $\gamma_1=\mathrm{EP}(X)=2(1-g)$, on trouve
\[
c_1=c_1^{\mathbb R}(E_0)=(1-g)
\]

\uline{N.B.} Le fibré vectoriel \inserted{réel} $E_0$ \deleted{(\unsure{on})}\inserted{sur $X$} \unsure{est} obtenu, d'une part par descente à partir d'un fibré vectoriel de rang $2$ \unsure{constt|const} sur $D$, à l'aide d'une repr. linéaire réelle (\uline{pas} \uline{complexe} !) de $\Gamma$, \deleted{\unsure{donc}}\inserted{donc} \deleted{\unsure{part, un fibré algébrique}} \inserted{à ce titre} $E_0$ est \unsure{muni} d'une \uline{connexion} \uline{intégrable} ; D'autre part, il a été obtenu aussi par descente : partant d'un fibré vectoriel \uline{complexe}\inserted{ (attention, pas holomorphe !) de rang $1$} sur $D$, à l'aide d'\unsure{opération|opérations} de $\Gamma$ sur ledit, (provenant en fait d'opérations \unsure{``de|de} $G$ sur le fibré, et définies par l'intermédiaire de la repr naturelle de $\widetilde S$), \deleted{\unsure{et}} d'où un fibré $E_0$ est muni d'une structure de fibré \uline{complexe} \uline{unitaire} de rang $1$. \unsure{Attention|Affirmation} que, bien entendu, la structure vectorielle \unsure{ou} complexe de ce fibré \uline{n'est pas} invariante par la connexion intégrable. Son carré tensoriel complexe \unsure{serait|serais} $T_X$ \deleted{\unsure{qui} \unsure{serait donc pourvu de connexion intégrable par}}

% margin note written sideways
\margin{\uline{Re. NB.} D'après la relation \unsure{analogue|marquée} pour les représentations de $\widetilde S$, on trouve la formule $E_0\otimes_{\mathbb R}E_0\simeq T_X+2\mathbb{1}_{\mathbb R}$ (Isom de fibrés vectoriels réels) \uline{Corollaire} $T_X+2\mathbb{1}_{\mathbb R}$ admet une connexion intégrable $\unsure{\pm|\ne}$ canonique. \unsure{cf. $(=)$ ci-dessus}}
\margin{Pour les structures réelles sous-jacentes \unsure{n'aura|n'aurai} pas de connexion intégrable en général (?)}
% archive sheet number: 8
\scanpage{11}
\mspage{\textcircled{2}}
% Reversed typescript shows through.
$\Gamma\subset G$ d'où
\diagram{\begin{tikzcd}[ampersand replacement=\&,column sep=small,row sep=small]
\& H^*(B_G,\mathbb Z)\arrow[dl]\arrow[d,"\unsure{\wr|2}"]\arrow[r,"\deleted{\unsure{\mathrm{res}}}"] \&
H^*(\deleted{\unsure{\Gamma}}B_{\Gamma},\mathbb Z)\arrow[d,phantom,"\rotatebox{90}{$\simeq$}" description]\\
H^*(B_{\widetilde G},\mathbb Z)\arrow[d,"\wr"] \&
H^*(B_S,\mathbb Z)\simeq\mathbb Z[\gamma_1]\arrow[dl]\arrow[r,bend right=36] \&
\begin{gathered}H^*(X,\mathbb Z)\\ \gamma_1^X\end{gathered}\\
H^*(B_{\widetilde S},\mathbb Z)\simeq\mathbb Z[c_1] \& \&
\end{tikzcd}}

\[
\deleted{\unsure{c|2c|\text{Car}}}\quad\gamma_1=2c_1
\]
d'où
\[
\boxed{\gamma_1^X\in H^2(X,\mathbb Z)\simeq\boxed{\mathbb Z}}
\]

À calculer $\gamma_1^X$ ??

\deleted{\unsure{Soit|Sat}} Quand on peut trouver $\widetilde\Gamma$,
\diagram{\begin{tikzcd}[ampersand replacement=\&]
B_\Gamma=X\arrow[rr]\arrow[dr,dashed] \& \& B_G\\
\& B_{\widetilde G}\arrow[ur] \&
\end{tikzcd}}
Alors on trouve $=$
\[
\boxed{\frac12\gamma_1^X=c_1^X\in H^2(X,\mathbb Z)=\mathbb Z}
\]
dont on peut donner une interprétation en termes \inserted{\unsure{de classe|des classes} $c_1^{\mathbb R}$} de fibré vectoriel réel $E$ sur $B_{\widetilde G}$, de rang $2$.

Le calcul donne $c_1=1-g$, donc
\[
\boxed{\gamma_1^X=2(1-g)=2-2g\deleted{\unsure{g|\{}}\unsure{=|\ne}\mathrm{EP}(X)\ne0}
\]
Dans le cas général, on \unsure{sait|voit} qu'il existe un sous-groupe d'indice fini $\Gamma'$ de $\Gamma$ qui se \unsure{remonte}.
\diagram{\begin{tikzcd}[ampersand replacement=\&,column sep=large]
B_{\Gamma'}=X'\arrow[r,"p","{\substack{\text{rev. étale}\\\text{de degré }d}}"' ] \& B_\Gamma=X\arrow[r] \& B_G
\end{tikzcd}}
On aura alors
\[
\gamma_1^{X'}=p^*(\gamma_1^X)\quad\text{\unsure{déf.|d'où|donc}}
\]
donc \unsure{vu qu'on est en}
\[
\gamma_1^{X'}=d\gamma_1^X
\]
Mais aussi, comme bien connu
\[
\mathrm{EP}(X')=d\,\mathrm{EP}(X)
\]
et comme $\gamma_1^{X'}=\mathrm{EP}(X')$, on conclut $\gamma_1^X=\mathrm{EP}(X)$ Ok.
% archive sheet number: 10
\scanpage{12}
\mspage{19}
% One continuous cancellation, made with a frame and diagonal strokes, reaches
% from (4.16) to the close of the paragraph before “Alors l'image”.
\[
\deleted{\handwritten{H^1(\unsure{\underline h|h|\underline k|k|y},\mathrm{Pic}_{X/h})=0}\quad
\deleted{\mathrm{Br}(X_{\handwritten{\unsure{y}}})=0,\quad H^{\handwritten3}(y,\underline G_m)=0,}\quad
\handwritten{\mathrm{Ker}(H^3(y,G_m)\longrightarrow H^3(X_y,G_m))}}
\tag*{\deleted{(4.1\deleted{\unsure{4}}\handwritten{6})}}
\]
\deleted{(Pour le voir, utiliser \deleted{\unsure{un|l'} homomorphisme de}\inserted{\handwritten{l'homomorphisme}} suite\handwritten{s} exacte\handwritten{s} du type (4.5) relatives à $\bar f:\bar X\handwritten{\longrightarrow}\bar Y$ \deleted{\unsure{et à $f_y:X_y\longrightarrow y$, qui donne,}}\inserted{\handwritten{et à $f_y:X_y\longrightarrow y$, qui s'écrit,}} \handwritten{compte tenu de (3.4) b) \unsure{:|;} et de (\phantom{000}) :}}
\diagram{\begin{tikzcd}[ampersand replacement=\&,column sep=small,row sep=normal,
execute at end picture={\draw (brNW.north west) rectangle (brSE.south east);
\draw (brNW.north east) -- (brSW.south west) (brNE.north east) -- (brSE.south west);}]
|[alias=brNW]| \deleted{\mathrm{Br}(\bar Y)}\arrow[r]\arrow[d,"\handwritten{\alpha}"] \&
|[alias=brNE]| \deleted{\mathrm{Br}(\bar X)}\arrow[r,"\handwritten{0}" description]\arrow[d,"\handwritten{\beta}"] \&
\deleted{H^1(\bar Y,\bar P)}\arrow[r]\arrow[d,"\handwritten{\gamma}"] \&
\deleted{H^{\handwritten3}(\bar Y,\underline G_m)}\arrow[r]\arrow[d,"\handwritten{\alpha'}"] \&
\deleted{H^{\handwritten3}(\bar X,\underline G_m)}\arrow[d,"\handwritten{\beta'}"]\\
|[alias=brSW]| \deleted{\mathrm{Br}(y)}\arrow[r] \& |[alias=brSE]| \deleted{\mathrm{Br}(X_y)}\arrow[r,"\handwritten{0}" description] \&
\deleted{H^1(y,P_{\unsure{o|0}})}\arrow[r] \& \deleted{H^{\handwritten3}(y,\underline G_m)}\arrow[r] \&
\deleted{H^{\handwritten3}(X_y,\underline G_m)}
\end{tikzcd}}
\deleted{\deleted{\handwritten{\unsure{?}}}\quad ).}

\deleted{Les flèches $\handwritten{\alpha,\beta,\alpha'}$ du diagramme sont bijectives, en vertu de (BR I\phantom{0000}), (\handwritten{3}.1), (\phantom{000}), d'où résulte que \deleted{il en est de même} $\handwritten{\gamma}:H^1(Y,P)\handwritten{\longrightarrow} H^1(\handwritten{y\deleted{\unsure{{}_y}}},P_o)$ est injective, où $P_o=\underline{\mathrm{Pic}}_{X_y/y}$. Donc (4.1\handwritten{5}) est une conséquence de $H^1(y,P_o)=0$, elle-\deleted{\unsure{m}}même conséquence d'après la deuxième suite exacte ci-dessus, des relations (4.1\handwritten{6})).}

Alors l'image de $H^1(Y,P)$ dans $H^1(Y,\underline B)$ est le sous-groupe de $H^1(Y,\underline B)$ formé des \unsure{élements|élemtns|éléments} dont l'image dans chaque $H^1(\bar Y,\underline B)$ est nulle. Ce sous-groupe s'appelle le \uline{groupe de Tate-Chafarévitch de} $Y$ \uline{à coéfficients dans} $\underline B$, et se note $\handwritten{\overline{\underline{\mathrm{III}}}}(Y,\underline B)$ (cf plus bas \inserted{pour} la relation avec la définition originelle de Tate et Chafarévitch). On trouve alors la suite exacte déduite de (4.1\deleted{\unsure{2|4}}\handwritten{4}) :
\[
\ldots\ \mathrm{Pic}(X_{\handwritten{\eta}}/\handwritten{\eta})\handwritten{\longrightarrow\coprod}_{y}H^1(y,\underline E_y)
\handwritten{\longrightarrow} H^1(Y,P)\handwritten{\longrightarrow\overline{\underline{\mathrm{III}}}}(Y,\underline B)\handwritten{\longrightarrow}0\ ,
\tag*{(4.1\deleted{\unsure{5|6}}\handwritten{7})}
\]
\margin{\handwritten{\textcircled{P}}}
explicitant $H^1(Y,\deleted{\unsure{\underline B}}\handwritten P)$ comme une extension du groupe de Tate-Chafarévitch par un quotient de la somme $\handwritten{\coprod}_yH^1(y,\underline E_y)$.

\margin{\handwritten{\unsure{PP|//}}}
Il faut enfin expliciter cette dernière somme. Notons que $\underline E_y$ \deleted{n'est non nul qu} est nul si $X_y$ est \inserted{\deleted{géométriquement}\handwritten{géométriquement}} intègre, à fortiori si $X_y$ est lisse sur $y$, de sorte que si $X_{\handwritten{\eta}}$ est séparable sur $k(\handwritten{\eta})$ (par exemple, lisse sur $k(\handwritten{\eta})$), alors le support de $E$ est un \deleted{ensemble}\inserted{contenu dans (en fait, égal à) l'ensemble} \uline{fini} des $y_i\handwritten{\in}Y$ tels que $X_{y_i}$ ne soit pas géométriquement intègre. Un calcul facile donne pour $\underline E_y$ la structure suivante (en procédant par descente à partir de $\underline E_{\bar y}$). On écrit \handwritten{l'égalité de \deleted{\unsure{cycles}}\inserted{diviseurs sur} $X$ :}
\[
X_y=\handwritten{\sum}_{i} a_y^i\,C_y^i\ ,
\tag*{(4;1\deleted{\unsure{6}}\handwritten8)}
\]
les $C_y^i$ étant les composantes irréductibles de $X_y$, et\inserted{\handwritten{ les}} $a_y^i$ leur multiplicité dans la fibre. On désigne par
% archive sheet number: 11
\scanpage{13}
\mspage{\textcircled{3}}
% Reversed typescript shows through.
\uline{NB} Le rev. universel de $G$ est de \deleted{\unsure{groupe}}\inserted{noyau} $\mathbb Z$. Il donne une extension \inserted{centrale} de $\Gamma$ par $\mathbb Z$, d'où un élément de $H^2(\Gamma,\mathbb Z)\simeq H^2(X,\mathbb Z)$. Montrons que c'est $\gamma_1$. Plus généralement, \unsure{quand on a|partant on a} une extension d'un gpe top. $G$ par un groupe discret $H$, lui correspond un elt de $H^2(B_G,H)$ (\unsure{fonctorielt|fonctoriel} en $G,H$ \ldots).
\margin{\diagram{\begin{tikzcd}[ampersand replacement=\&,column sep=tiny]
B_H\arrow[r]\&B_E\arrow[r]\&B_G
\end{tikzcd}}}
Dans le cas où $G$ est $S=$ tore ``\unsure{standart|standard}, et $\widetilde S=$ rev. universel de $\widetilde S$, de noyau $\mathbb Z$, l'\unsure{élém.} qu'on \unsure{trouve} \unsure{dans} $H^{\unsure{2|*}}(B_S,\mathbb Z)$ n'est autre que le générateur canonique de Chern.

Ces remarques ensemble \unsure{établissent|établ\textsuperscript{t}} \unsure{expmt} l'interprétation \unsure{signalée|esquissée} de $\gamma_1^X\in$ \deleted{\unsure{dans}} $H^2(X,\mathbb Z)\simeq H^2(\Gamma,\mathbb Z)$. Donc son image dans $H^2(\Gamma,\mathbb Z/2\mathbb Z)$ \deleted{\unsure{est nulle}}\inserted{$=\mathbb Z/2\mathbb Z$} est la classe de l'extension de $\Gamma$ par $\mathbb Z/2\mathbb Z$, déduite de $\widetilde G\longrightarrow G$. \unsure{C'est nulle|C'est celle} donc \deleted{\unsure{ssi}} \unsure{précisément} \unsure{si scindée|si elle}. Comme $\mathrm{EP}(X)=2(1-g)$ est pair, elle est tjrs nulle !

On gagne ! Il existe tjrs un $\widetilde\Gamma$ !

\textcircled{4} Revenant, avec cette information, à la situation \textcircled{1}, on voit que \deleted{\unsure{la}}, pour tte courbe alg. \unsure{(non complète|(non complétée|lisse complète} connexe de genre $g\geq2$, fibré tangent $T_X$, le fibré vectoriel \deleted{complexe}\inserted{réel de rang $\unsure{4|2|1}$} $T_X\deleted{\unsure{\otimes_{\mathbb R}T_X}}+2\mathbb{1}_{\mathbb R}$ sur $X$ est muni d'une connexion intégrable, (attention, elle ne respecte pas les
% archive sheet number: 12
\scanpage{16}
% Mirror-image typescript on the other side.
$G$ gpe de Lie compact connexe

\diagram{\begin{tikzcd}[ampersand replacement=\&,column sep=small,row sep=small]
C_{(0)}^{p,q}\arrow[d,phantom,"\cup" description] \& {}=\Omega^q(\unsure{\mathord{-}}G^{p+1}/G) \&\\
C_{(1)}^{pq}\arrow[d,phantom,"\cup" description] \& {}=\omega^{(q)}(G^{p+1}/G) \&\\
C_{(2)}^{pq} \& {}=\omega_{\mathrm{inv}}^{(q)}(G^{p+1}/G) \&
\end{tikzcd}}
\margin{inclusion induit un isom sur \unsure{les termes|le terme|les dernières} $\,{}'E_1$ (cohomologie pour $p$ fixe}
\margin{id}

Dans ce dernier bicomplexe, l'opérateur différentiel par rapport à $q$ \uline{est nul}. Donc on trouve
\[
\begin{aligned}
H^*(B_G,\mathbb R)&=H^*(C_{\bullet}^{pq})\simeq H^*(C_{(1)}^{pq})
\simeq\sum_{p,q}\bigl\{H^p\bigl(\unsure{\nu|v}\longrightarrow
\omega_{\mathrm{inv}}^{(q)}(G^{\nu+1}/\unsure{G})\unsure{)\}|)}\\[-2pt]
&\hspace{15em}\rotatebox{90}{$=$}\\
&\hspace{15em}E_2^{pq}
\end{aligned}
\]
(isomorphismes canoniques)

\uuline{La suite spectrale dégénère} : différentielles $d_r$ nulles pour $r\geq2$.

Soit $T\subset\omega^*_{\mathrm{inv}}(G)$, \unsure{élémts|éléments} transgressifs. (\unsure{Ils|OK} sont de degrés\unsure{$+$|s} impairs)

On a
\[
T\xrightarrow{\varphi} E_2^{1*}=\omega_{\mathrm{inv}}^*(G\times G/G)
\subset\underset{\substack{\rotatebox{90}{$=$}\\
  \omega_{\mathrm{inv}}^*(G)\otimes\omega_{\mathrm{inv}}^*(G)}}
{\omega_{\mathrm{inv}}^*(G\times G)}
\]
par
\[
\varphi(T)=T\otimes1-1\otimes T
\]
\uline{Th} . $\mathrm{Sym}^*(T)\longrightarrow H^*(B_G,\mathbb R)$ \uline{déduit} de $\varphi$ \uline{est} \uline{un isomorphisme}.
% archive sheet number: 15
\scanpage{17}
% Three long oblique pen strokes cancel the typed passage as a whole.
\deleted{triangulée $\deleted{\unsure{a|c}}C(X,A)$ associée ainsi à $X$ détermine l'ensemble des composantes connexes de $X$, et le groupe fondamental de chacune des composantes connexes, et bien entendu la cohomologie de chacune, à coéffieients des $A$-modules tordus de type fini \ldots{} Lorsque $A=\underline Z$, il semble bien que Hurewicz (ou une variante de H. mise au point par Artin-Mazur) montre qu'un morphisme $f:X\handwritten{\to}Y$ est un isomorphisme dans la catégorie homotopique, si et seulement si le foncteur $C(Y,\underline Z)\handwritten{\to}C(\mathbb X,\underline{\mathbb Z})$ est une équivalence. On se demande alors si on ne peut pas réconstituer entièrement $X$ en termes de $C(X,\underline Z)$, du moins (si on a fait dans la définition de $C(X,\underline Z)$ les\handwritten{/}hypothèses de finitude que j'ai dites) si on se restreint à des $X$ satisfaisant certaines conditions de finitude (comme des polyèdres qui sont finis en chaque dimension). Plus précisément, serait-il vrai que tout f\unsure{a|o}ncteur exact $C(Y,\underline Z)\handwritten{\to}C(X,\underline Z)$ proviendrait à isomorphisme près d'un unique morphisme $X\handwritten{\to}Y$ ? Même si ceci était trop optimiste, les catégories $C(X,A)$ me semblent des invariants rupinants, qui doivent être des approximations très bonnes de l'objet $X$ à homotopie près, et on aimerait expliciter ce point de vue en \deleted{\unsure{explicitant|expliciter}} interprétant en \deleted{\unsure{temmes|termes}}termes de ces catégories les constructions les plus importantes familières aux homotopistes: construction des groupes d'homotopie supérieurs, plus \deleted{\unsure{génralement|généralement}}\inserted{précisément} des espaces de lacets, et en sens inverse des espaces classifiants, ou des suspensions, produits cartésiens d'espaces, \inserted{\handwritten{squelettes, \unsure{tuages|tuer|étages|images} des groupes d'homotopie supérieurs}} etc. \deleted{Notes aussi que \unsure{no|ne}} C'est là un complexe de questions manifestement très lié à celui que j'a\uline{i} abordé dans ma lettre du 26.3. dont tu viens de m'envoyer une photocopie, et je crois que quelqu'un devrait bien un jour regarder ces choses systématiquement. Notes que si tu fais sur $X$ une hypothèse de finitude comme plus haut, alors $C(X,\underline{\unsure{\mathbb Z|Z}})$ sera}
% archive sheet number: 16
\scanpage{18}
$G$ groupe de Lie \deleted{compact}\inserted{complexe réductif} connexe. Considérons les conditions ($p\in\mathbb P$ donné)

\[
\begin{aligned}
\mathrm{(i)}\quad &G\text{ sans }p\text{-torsion}\\
\mathrm{(ii)}\quad &B_G\text{ sans }p\text{-torsion}\\
\mathrm{(iii)}\quad &\operatorname{Sym}^{*}_{\mathbb Z}(\widehat B)\longrightarrow H^*(G/B,\mathbb Z)
  \quad\text{a une image d'indice}\\
&\text{premier à }p,\quad\text{i.e.}\quad
  \operatorname{Sym}_{\deleted{\unsure{\mathbb Z}}\mathbb F_p}(\widehat B\otimes_{\mathbb Z}\mathbb F_p)\to H^{\unsure{*|\circ|2}}(G/B,\mathbb F_p)\\
&\text{est surjectif.}
\end{aligned}
\]

\[
\begin{aligned}
\mathrm{(i\ bis)}\quad &{\unsure{H|\mathbb H}}^{*}(G,\mathbb F_p)\simeq\bigwedge^{*}V,\quad V\text{ vectoriel gradué \inserted{de rang $r$} sur }\mathbb F_p,\\
&\text{formé des éléments primitifs.}\\
\mathrm{(ii\ bis)}\quad &H^{\unsure{\circ|*}}(B_G,\mathbb F_p)\text{ est une algèbre de polynômes}\\
&\text{sur }\mathbb F_p\text{ à }r\text{ générateurs.}
\end{aligned}
\]

\par\noindent\hfil\rule{.47\linewidth}{.4pt}\hfil\par

\deleted{(i) N.B. Dans (i bis), on doit avoir}

$\mathrm{(i)}\deleted{\unsure{\Longleftarrow}}\Longrightarrow\mathrm{(i\ bis)}$ par th. de Hopf. (cf Notes de Borel\ldots{})

$\mathrm{(i\ bis)}\Longrightarrow\mathrm{(ii\ bis)}$ par suite spectrale \deleted{\unsure{de|de la}} descente, ou th. de transgression de Borel.

$\mathrm{(ii\ bis)}\Longleftrightarrow\mathrm{(iii)}$ [Suite spectrale de Leray dans fibration $G/B\unsure{\leftarrow|\hookleftarrow}\mathbb B_B\mathbin{\unsure{-|\to}}B_G$ est triviale] \deleted{\unsure{de}}

$\mathrm{(ii\ bis)}\Longrightarrow\mathrm{(ii)}$ vu égalité des nbs de Betti, coeff $\mathbb Q$ ou $\mathbb F_p$

\diagram{\begin{tikzcd}[ampersand replacement=\&,column sep=small,row sep=small]
\mathrm{(i)}\arrow[r,Rightarrow]\arrow[d,Rightarrow]\&\mathrm{(i\ bis)}\arrow[dl,Rightarrow]\\
\mathrm{(ii\ bis)}\arrow[r,Leftrightarrow]\arrow[d,Rightarrow]\&\mathrm{(iii)}\\
\mathrm{(ii)}\&
\end{tikzcd}}
% archive sheet number: 17
